% Generated draft: requires mathematical and source review.
\documentclass[11pt]{article}
\usepackage[margin=25mm]{geometry}
\usepackage{fontspec}
\setmainfont{Latin Modern Roman}
\usepackage{amsmath,amssymb,mathtools,unicode-math}
\setmathfont{Latin Modern Math}
\usepackage{enumitem,needspace,xurl,hyperref}
\setlength{\parindent}{0pt}
\setlength{\parskip}{5pt}
\setlength{\emergencystretch}{3em}
\newcommand{\N}{\mathbb{N}}
\newcommand{\Z}{\mathbb{Z}}
\newcommand{\Q}{\mathbb{Q}}
\newcommand{\R}{\mathbb{R}}
\newcommand{\C}{\mathbb{C}}
\newcommand{\F}{\mathbb{F}}
\newcommand{\A}{\mathbb{A}}
\newcommand{\PP}{\mathbb{P}}
\newcommand{\E}{\mathbb{E}}
\newcommand{\eps}{\varepsilon}
\newcommand{\dd}{\,\mathrm{d}}
\begin{document}
\section*{UnsolvedMath ID: 30000725}
\textbf{Source problem code:} catalogue entry
\subsection{Definitions and mathematical statement}
For a finite simple undirected graph \(G=(V,E)\), the standard Colin de Verdi\`ere number \(\mu(G)\) is the maximum nullity of a real symmetric matrix \(M\), indexed by \(V\), satisfying: \(M_{uv}<0\) for \(uv\in E\), \(M_{uv}=0\) for distinct nonadjacent \(u,v\), exactly one negative eigenvalue, counted with multiplicity, and the strong Arnold property. The latter means that the only symmetric matrix \(X\) satisfying \(MX=0\) and \(X_{uv}=0\) whenever \(u=v\) or \(uv\in E\) is \(X=0\).

Form the panelling \(\widetilde G\) by attaching one closed two-dimensional disk along each circuit of \(G\); under the usual simple-cycle convention the attaching map identifies its boundary with that cycle. Write \(\sigma(G)\) for the even-mapping graph invariant named in the source. Its defining class of complexes and maps and its parity convention are not specified in the supplied record, so this symbol requires the definition from the cited report or paper; it must not be replaced merely by an embedding dimension of the two-dimensional panelling.

The extracted question is whether, for every undirected graph \(G\) in the intended domain,
\[
 \mu(G)\leq5
 \quad\Longleftrightarrow\quad \sigma(G)\leq5
 \quad\Longleftrightarrow\quad
 \widetilde G\text{ embeds in }\R^4.
\]
In the finite simple topological interpretation, the last condition means there exists a continuous map \(e:\widetilde G\to\R^4\) that is a homeomorphism onto its image. The record does not settle whether the intended category is topological or piecewise linear, nor whether graphs beyond the finite simple case are included.

The supplied assessment, dated 2026-08-21, reports \(\mu(G)\leq\sigma(G)\), which yields \(\sigma(G)\leq5\Rightarrow\mu(G)\leq5\), and reports strict gaps between these parameters. It says that the full extracted threshold-five equivalence is not established. A strict gap alone does not establish a failure of this particular threshold equivalence.
\subsection{Short English statement}
The question compares three proposed characterizations at the threshold five: a matrix invariant, an even-mapping invariant, and embeddability of the complex obtained by filling every circuit in four-dimensional Euclidean space. The supplied later assessment gives one implication and reports partial progress; exact graph and mapping conventions still require source review.
\subsection{Sources}
\begin{itemize}
\item Local UnsolvedMath registry, permanent ID \(30000725\), exact source code \texttt{catalogue entry}, proposed year 2007; no individual proposer is supplied.
\item \emph{Graph Theory (2008)}, Oberwolfach Reports, DOI 10.4171/owr/2007/16, pp. 887--944. \url{https://doi.org/10.4171/owr/2007/16}.
\item Supplied literature source labelled ``even maps and graph parameters'': \url{https://arxiv.org/abs/1907.05055}.
\item The registry marks the statement \texttt{corrected\_verified}, explaining that corrupted list delimiters and geometric notation were repaired. Its \texttt{partially\_solved} label and assessment checked 2026-08-21 are supplied dated literature triage, requiring independent verification of status and cited claims.
\end{itemize}
\end{document}
